% -*- coding: utf-8 -*-
% !TEX program = xelatex
\documentclass[plain,noanswer,medmath]{../../template/jnuexam}
\usepackage{../../template/kysx}

\begin{document}


\examtitle{1997}{3}


%\loadproblems{1}{1997P1}%载入试卷一


\makepart{填空题}{本题共5分，每小题3分，满分15分}


\begin{problem}
设$y = f(\ln x) \e^{f(x)}$，其中$f$可微，则$\dy =$ \fillin{}．
\end{problem}


\begin{problem}
若$f(x) = \frac{1}{1 + x^2} + \sqrt{1 - x^2} \int_0^1f(x)\dx$，则$\int_0^1f(x)\dx =$ \fillin{}．
\end{problem}


\begin{problem}
差分方程$y_{t + 1} - y_t = t2^t$的通解为 \fillin{}．
\end{problem}


\begin{problem}
若二次型$f(x_1,x_2,x_3) = 2x_1^2+ x_2^2+ x_3^2+ 2x_1x_2 + tx_2x_3$是正定的，则$t$的取值范围是 \fillin{}．
\end{problem}


\begin{problem}
设随机变量$X$和$Y$相互独立且都服从正态分布$N(0,3^2)$，
而$X_1, \cdots ,X_9$和$Y_1, \cdots ,Y_9$分别是来自总体$X$和$Y$的简单随机样本，
则统计量$U = \frac{X_1 + \cdots + X_9}{\sqrt{Y_1^2+ \cdots + Y_9^2}}$服从分布\fillin{}，参数为 \fillin{}．
\end{problem}


\makepart{选择题}{本题共5小题，每小题3分，满分15分}


\begin{problem}
设$f(x) = \int_0^{1 - \cos x} \sin t^2 \dt$, $g(x) = \frac{x^5}{5} + \frac{x^6}{6}$，
则当$x \to 0$时，$f(x)$是$g(x)$的\pickout{}
\begin{abcd}
\item 低阶无穷小．                    
\item 高阶无穷小．
\item 等价无穷小．                    
\item 同阶但不等价的无穷小．
\end{abcd}
\end{problem}


\begin{problem}
若$f(- x) = f(x)$（$-\infty < x < +\infty$），在$(-\infty ,0)$内$f'(x) > 0$，且$f''(x) < 0$，则在$(0, +\infty)$内有\pickout{}
\begin{abcd}
\item $f'(x) > 0$,$f''(x) < 0$．
\item $f'(x) > 0$,$f''(x) > 0$．
\item $f'(x) < 0$,$f''(x) < 0$．
\item $f'(x) < 0$,$f''(x) > 0$．
\end{abcd}
\end{problem}


\begin{problem}
设向量组$\alpha_1$,$\alpha_2$,$\alpha_3$线性无关，则下列向量组中，线性无关的是\pickout{}
\begin{abcd}
\item $\alpha_1 + \alpha_2$, $\alpha_2 + \alpha_3$, $\alpha_3 - \alpha_1$．
\item $\alpha_1 + \alpha_2$, $\alpha_2 + \alpha_3$, $\alpha_1 + 2\alpha_2 + \alpha_3$．
\item $\alpha_1 + 2\alpha_2$, $2\alpha_2 + 3\alpha_3$, $3\alpha_3 + \alpha_1$．
\item $\alpha_1 + \alpha_2 + \alpha_3$, $2\alpha_1 - 3\alpha_2 + 22\alpha_3$,
      $3\alpha_1 + 5\alpha_2 - 5\alpha_3$．
\end{abcd}
\end{problem}


\begin{problem}
设$A,B$为同阶可逆矩阵，则\pickout{}
\begin{abcd}
\item $AB = BA$．
\item 存在可逆矩阵$P$，使$P^{-1} AP = B$．
\item 存在可逆矩阵$C$，使$C^T AC = B$．
\item 存在可逆矩阵$P$和$Q$，使$PAQ = B$．
\end{abcd}
\end{problem}


\begin{problem}
设两个随机变量$X$与$Y$相互独立且同分布：
$$P\{X = - 1\} = P\{Y = - 1\} = \frac{1}{2},\qquad P\{X = 1\} = P\{Y = 1\} = \frac{1}{2},$$
则下列各式中成立的是\pickout{}
\begin{abcd}
\item $P\{X = Y\} = \frac{1}{2}$．
\item $P\{X = Y\} = 1$．
\item $P\left\{X + Y = 0 \right\} = \frac{1}{4}$．
\item $P\{XY = 1\} = \frac{1}{4}$．
\end{abcd}
\end{problem}


\makepart{}{本题满分6分}

\begin{problem}
在经济学中，称函数$Q(x) = A[\delta K^{- x} + (1 - \delta)L^{- x}]^{- \frac{1}{x}}$为固定替代弹性生产函数，
而称函数$\overline{Q}= AK^{\delta}L^{1 - \delta}$为Cobb-Douglas生产函数(简称C--D生产函数).
试证明：当$x \to 0$时，固定替代弹性生产函数变为C--D生产函数，即有$\lim_{x \to 0} Q(x) = \overline{Q}$．
\end{problem}


\makepart{}{本题满分5分}

\begin{problem}
设$u = f(x,y,z)$有连续偏导数，$y = y(x)$和$z = z(x)$分别由方程
$\e^{xy} - y = 0$和$\e^x - xz = 0$所确定，求$\frac{\du}{\dx}$．
\end{problem}


\makepart{}{本题满分6分}

\begin{problem}
一商家销售某种商品的价格满足关系$p = 7 - 0.2x$(万元/吨),
$x$为销售量(单位：吨)，商品的成本函数$C = 3x + 1$(万元).
\begin{enumerate}
\item 若每销售一吨商品，政府要征税$t$(万元)，求该商家获最大利润时的销售量；
\item $t$为何值时，政府税收总额最大．
\end{enumerate}
\end{problem}


\makepart{}{本题满分6分}

\begin{problem}
设函数$f(x)$在$[0, +\infty)$上连续、单调不减且$f(0) \ge 0$，试证函数
$$F(x) = \begin{cases}
  \frac{1}{x}\int_0^x t^n f(t)\dt , & \text{若}x > 0, \\
                                 0, & \text{若}x = 0,
\end{cases}$$
在$[0, +\infty)$上连续且单调不减(其中$n > 0$).
\end{problem}


\makepart{}{本题满分6分}

\begin{problem}
从点$P_1(1,0)$作$x$轴的垂线，交抛物线$y = x^2$于点$Q_1(1,1)$；
再从$Q_1$作这条抛物线的切线与$x$轴交于$P_2$，然后又从$P_2$作$x$轴的垂线，交抛物线于点$Q_2$,
依次重复上述过程得到一系列的点$P_1,Q_1;P_2,Q_2; \cdots ;P_n,Q_n; \cdots$．\par
\begin{enumerate*}
\item 求$\overline{OP_n}$；
\item 求级数$\overline{Q_1P_1} + \overline{Q_2P_2} + \cdots + \overline{Q_n P_n} + \cdots$的和．
\end{enumerate*}\par
其中$n(n \ge 1)$为自然数，而$\overline{M_1M_2}$表示点$M_1$与$M_2$之间的距离．
\end{problem}


\makepart{}{本题满分6分}

\begin{problem}
设函数$f(t)$在$[0, +\infty)$上连续，且满足方程
$$f(t) = \e^{4\pi t^2} + \iint_{x^2 + y^2 \le 4t^2} f\left(\frac{1}{2}\sqrt{x^2 + y^2}\right) \dx\dy,$$
求$f(t)$．
\end{problem}


\makepart{}{本题满分6分}

\begin{problem}
设$A$为$n$阶非奇异矩阵，$\alpha$为$n$维列向量，$b$为常数．记分块矩阵
$$P = \begin{pmatrix}
              E & 0 \\
  -\alpha^T A^* & |A|
\end{pmatrix},\quad
Q = \begin{pmatrix}
         A & \alpha \\
  \alpha^T & b
\end{pmatrix},$$
其中$A^*$是矩阵$A$的伴随矩阵，$E$为$n$阶单位矩阵．
\begin{enumerate}
\item 计算并化简$PQ$；
\item 证明：矩阵$Q$可逆的充分必要条件是$\alpha^T A^{-1}\alpha \ne b$．
\end{enumerate}
\end{problem}


\makepart{}{本题满分10分}

\begin{problem}
设三阶实对称矩阵$A$的特征值是$1$, $2$ ,$3$；矩阵$A$的属于特征值$1$, $2$的特征向量分别是
$\alpha_1 = (- 1, - 1,1)^T,\alpha_2 = (1, - 2, - 1)^T$．\par
\begin{enumerate*}
\item 求$A$的属于特征值$3$的特征向量；
\item 求矩阵$A$．
\end{enumerate*}
\end{problem}


\makepart{}{本题满分7分}

\begin{problem}
假设随机变量$X$的绝对值不大于$1$；$P\{X = - 1\} = \frac{1}{8},P\{X = 1\} = \frac{1}{4}$；
在事件$\{- 1 < X < 1\}$出现的条件下，$X$在$(- 1,1)$内的任一子区间上取值的条件概率与该子区间长度成正比．
试求$X$的分布函数$F(x) = P\{X \le x\}$．
\end{problem}


\makepart{}{本题满分6分}

\begin{problem}
游客乘电梯从底层到电视塔顶层观光；电梯于每个整点的第$5$分钟、$25$分钟和$55$分钟从底层起行．
假设一游客在早晨八点的第$X$分钟到达底层候梯处，且$X$在$[0,60]$上均匀分布，求该游客等候时间的数学期望．
\end{problem}


\makepart{}{本题满分6分}

\begin{problem}
两台同样自动记录仪，每台无故障工作的时间服从参数为$5$的指数分布；
首先开动其中一台，当其发生故障时停用而另一台自行开动．
试求两台记录仪无故障工作的总时间$T$的概率密度$f(t)$、数学期望和方差．
\end{problem}


\end{document}
